\thispagestyle{empty}
\begin{tcolorbox}[enhanced, colback=hblue, colframe=hblue, boxrule=0pt, arc=4pt, left=16pt, right=16pt, top=22pt, bottom=22pt]
  \color{white}\sffamily
  {\large\bfseries Bewijzen in de Wiskunde \textperiodcentered\ 2026--2027}\\[10pt]
  {\fontsize{34}{38}\selectfont\bfseries Workbook}\\[6pt]
  {\fontsize{20}{24}\selectfont Chapters 0--5}\\[14pt]
  {\normalsize Minimal theory, taken from the book's own definitions, strategies and results, followed immediately by the book's own exercises: every exercise of Chapters 0--5 (sections 0.1--5.2 and the chapter exercises 0.E--5.E) is either worked in the text or solved in Appendix S.}
\end{tcolorbox}

\vspace{10pt}
\begin{tcolorbox}[enhanced, colback=white, colframe=hblue2, boxrule=0.8pt, arc=2pt, left=10pt, right=10pt, top=8pt, bottom=8pt]
\sffamily\small\sloppy
\begin{itemize}[leftmargin=1.2em]
\item \textbf{\textcolor{hblue}{Source.}} C.~Newstead, \emph{An Infinite Descent into Pure Mathematics}, Utrecht adaptation (version of 29 July 2026), pages 1--180. All numbers (Definition 0.36, Strategy 1.1.28, Exercise 3.2.40, Theorem 5.2.26, \dots) are the book's numbers; nothing here comes from outside the book except the standard algebra and calculus the book itself takes for granted.
\item \textbf{\textcolor{hblue}{Which exercises are worked where.}} Exercises flagged in the course planning (\fE\ essential, \fR\ recommended, bonus) are worked in full in the chapter text. All other book exercises (\fO) appear in NOW YOU boxes and are solved in Appendix S, so that you can attempt \emph{any} question in the book and check your work.
\item \textbf{\textcolor{hblue}{How to use it.}} Read one theory box, cover the solution of the worked exercise that follows, attempt it, then compare. Do the NOW YOU exercises before looking at Appendix S. The one-page toolbox in Appendix T lists the first line of every kind of proof.
\end{itemize}
\end{tcolorbox}

\vspace{8pt}
\hsub{How this workbook works}
Every section is a sequence of short cycles: \textbf{one piece of theory, then the exercises that use it.}

\begin{center}\sffamily\small
\begin{tabular}{@{}p{0.30\linewidth}p{0.65\linewidth}@{}}
\colorbox{hyellow}{\strut\textbf{Definition n.m}} & The book's definition, condensed. Learn it word for word: every proof starts by unfolding one of these.\\[3pt]
\colorbox{hlight}{\strut\textbf{Theorem / Proposition / Axiom}} & A result you may quote by number in a proof, with the idea of the book's proof.\\[3pt]
\colorbox{hgreen}{\strut\textbf{Strategy n.m}} & The book's proof strategies. Each one is a rule for what to write next.\\[3pt]
\colorbox{white}{\strut\textcolor{hblue}{\textbf{Exercise n.m}} \textcolor{horange}{\small\textbf{badge}}} & A book exercise worked in full, in the style the book uses; the badge gives its planning status. Cover the solution and try first.\\[3pt]
\colorbox{hpurplebg}{\strut\textbf{NOW YOU}} & The remaining book exercises for that piece of theory. Do them; every one is solved in Appendix S.\\[3pt]
\colorbox{hredbg}{\strut\textbf{Watch out}} & The book's own warnings (``Common error'', asides) and the places where marks are lost.
\end{tabular}
\end{center}

\begin{key}[{Flags from the course planning}]
\fE\ = \textbf{essential} exercise in the planning \quad \fR\ = \textbf{recommended} \quad \fO\ = in the book but not in the planning (still solved here). Section 4.1 (Peano's axioms) is not in the planning and is covered briefly; the polynomial and complex-number part of Chapter 0 is outside the course but its exercises are solved.\\[3pt]
Quiz 1 (23 Sep): numbers and \S1.1. \quad Quiz 2 (7 Oct): sets and functions. \quad Quiz 3 (21 Oct): relations and modular arithmetic. \quad Exam 4 Nov.
\end{key}

\begin{key}[{Conventions of this book}]
\begin{itemize}
\item $\N=\{0,1,2,\dots\}$: $0$ is a natural number (Definition 0.6). $[n]=\{k\in\N\mid k<n\}$ (Definition 0.25), so $[0]=\varnothing$ and $[3]=\{0,1,2\}$.
\item Truth values in truth tables are $\checkmark$ (true) and $\times$ (false), as in the book.
\item ``$b$ divides $a$'' means $a=qb$ for some $q\in\Z$ (Definition 0.36); it is about multiplication, never about the fraction $a/b$.
\item ``Odd'' means ``not even'' (Definition 0.41); that $a$ is odd iff $a=2b+1$ is a \emph{proposition} (1.1.58), not the definition.
\item $f[U]$ and $f^{-1}[V]$ are images and preimages (Definitions 3.1.31, 3.1.37); $[a]_\sim$ and $[a]_n$ are equivalence and congruence classes (Definitions 5.2.12, 5.2.18).
\item $\square$ ends a proof. \reason{grey text in braces} names the definition, strategy or rule used for a step.
\end{itemize}
\end{key}
